\documentclass{beamer} 


%Theme (more on this next week)
\usetheme{AnnArbor}

% Preamble:

%Here's where packages go. 
\usepackage{amsmath} %includes lots of mathematical symbols and content
\usepackage{lipsum} % For blind text
%\usepackage{fullpage}
%\usepackage[margin=1in]{geometry}
%\usepackage{geometry}
%\geometry{
%  top=1in,
%  inner=1in,
%  outer=1in,
%  bottom=1in,
%}
\usepackage{enumerate}
\usepackage{array,tabularx,multirow} %Helpful for certain tables and table features
\usepackage{graphicx}
\graphicspath{ {./Images/} }
\usepackage{wrapfig}
\usepackage{animate}
\usepackage{tikz}
\usetikzlibrary{arrows}
%\usepackage[all]{xypic}
\usepackage{xcolor}
% \usepackage{minted}
\usepackage{hyperref}%%Generally hyperref should be the last package loaded. 
%\hypersetup{
%    colorlinks=true,
%    urlcolor=blue,
%    linkcolor=black
%}


% To make typesetting augmented matrices easier. Code taken from
% https://tex.stackexchange.com/questions/2233/whats-the-best-way-make-an-augmented-coefficient-matrix 
% \makeatletter
% \renewcommand*\env@matrix[1][*\c@MaxMatrixCols c]{%
%   \hskip -\arraycolsep
%   \let\@ifnextchar\new@ifnextchar
%   \array{#1}}
% \makeatother

%%Custom Commands
% \newcommand{\Z}{\mathbb{Z}}
% \newcommand{\twobytwo}[4]{\begin{bmatrix} #1 & #2 \\ #3 & #4 \end{bmatrix}}
% \DeclareMathOperator{\Ker}{Ker}
% \DeclareMathOperator{\coker}{coker}

\title{Math 3680: Week 12}
\author{Arjan Singh Ellingson}
%\date{August 24th, 2026}
\date{\today}
\institute[Cal Poly]{California Polytechnic State University}
\begin{document}

%%%%%%%%%%%%%%%%%%%%%%%%%%%%%%%%


\begin{frame}
    \maketitle
\end{frame}
% I am copying an existing presentation I made earlier in october in an SE class on GraphQL methods. Hopefully this isn't too boring of a read...
% i also prefer using minted for code blocks over verbatim and texttt but it is a janky install

\begin{frame}
    \frametitle{Introduction}
    GraphQL (you may see gql) may is an API method that allows for communication between a client or user and a server.

    GraphQL is an alternative to the more common \texttt{REST} format used in web development.

    Calls are executed from files with the extension \texttt{.gql} or \texttt{.graphql}.

    It is not limited to a specific database or backend language.
\end{frame}

\begin{frame}
    \frametitle{Brief History \& Alternatives}
    \begin{enumerate}
        % these should make three slides shown sequentially 
        \item<1->started by Facebook in 2012, aiming to becoming open source by 2015
        \item<2->2018: GraphQL was adopted by GraphQL Foundation, with parent Linux Foundation\footnote{https://en.wikipedia.org/wiki/GraphQL}

        \item<3->Currently: APIs of companies like Facebook, GitHub (\texttt{gh cli}), Yelp, and Shopify utilize gql format
        \item<4->GraphQL combines each of REST's endpoints into one endpoint (usually called \texttt{/graphql})
        \item<5->The Problem: Rest often \textbf{overfetches} or \textbf{underfetches} the data required.
    \end{enumerate}
\end{frame}

% requires fragile because call is in a verbatim env
\begin{frame}[fragile]
    \frametitle{Structure of a GraphQL Call}
    A basic call is formatted as
    \begin{verbatim}     
schema {     
    query: Query     
    mutation: Mutation     
}     
type Query {}     
type Mutation {}     
\end{verbatim}
\end{frame}


\begin{frame}[fragile]     
\frametitle{Example from a different Coding Project}  

\tiny
\begin{columns}[T,onlytextwidth]
\begin{column}{0.33\textwidth}
\begin{semiverbatim}
\pause
    type ArtistNode {
      id: ID!
      label: String!
      genre: String
      era: String
      image: String
      kind: NodeKind
      visible: Boolean!
    }
\pause
    type RelationshipEdge {
      id: ID!
      source: ArtistNode!
      target: ArtistNode!
      type: String!
      directed: Boolean!
      year: Int
      note: String
      song: SongRef
    }
\end{semiverbatim}
\end{column}
\begin{column}{0.66\textwidth}
\begin{semiverbatim}
\pause
    type Query {
      artists(visibleOnly: Boolean = false): [ArtistNode!]!
    
      artist(id: ID!): ArtistNode
    
      edges(type: String, targetLabel: String): [RelationshipEdge!]!
    }
\end{semiverbatim}
\end{column}
\end{columns}

\end{frame}

\end{document}



